% discussion.tex -- Discussion
\section{Discussion}

Changes in investor decision rules generated market instability, while
available trading resources determined whether the shift toward selling
accompanied a small decline or a price crash. When enough traders received a
new rule for responding to their neighbors, a majority entered Sell even
though the remaining traders kept their original rules. The new rule
remained in place when a selected trader stopped selling, so the changed
response applied whenever the same combination of neighboring decisions
occurred again. Asking those traders to sell once changed only their current
decision and produced no majority in the 50 runs tested for each condition.
Keeping Rule~19286 in use produced a majority in every run at the two larger
fractions, and several other changes to a single response also produced
majorities, although most rules tested did not. Widespread selling therefore
depended on which response changed as well as on how many traders received
the new rule. Comparing markets with identical histories isolated the effect
of the rule change. Repeating that comparison across starting markets then
measured how often a majority formed at each fraction of traders given the
new rule.

Market instability could develop through neighboring traders' decisions
without traders responding to a price change. Keeping the technical traders in
Hold removed the influence of price on decisions inside the grid, yet Sell
majorities still developed. In the displayed runs that reached a majority,
groups of sellers formed, enlarged, and joined together. Local decisions
thus produced growing spatial patterns, as occurs in other models of
interacting sites away from
equilibrium~\cite{hinrichsen2000nonequilibrium,bray1994theory}. The
dependence on the fraction of traders given the new rule resembled models in
which the size of an initially changed group determines whether a change
spreads across a
population~\cite{granovetter1978threshold,watts2002simple,aizenman1988metastability,holroyd2003sharp}.
The fraction needed for a majority depended on both the number of
interacting traders and the time allowed. On larger grids, the probability
rose from low to high over a narrower range of rule adoption, but allowing
more time continued to lower the fraction at which half the runs reached a
majority. Extending the displayed markets also showed that a market without
a Sell majority at the original endpoint could later develop one alongside
a large price decline (Fig.~\ref{fig:dose-panels}). The persistence of mixed
decisions over a shorter period therefore did not establish lasting price
stability. We did not identify a threshold that would remain fixed
in larger markets or longer runs. Establishing such a limit would require
the estimates to stop changing with observation time and a mathematical
model for their dependence on grid
size~\cite{binder1981finite,landau2021guide}. The growing groups suggested an analogy with droplets forming in a
phase change, but their appearance did not establish the physical
mechanism of classical nucleation. We did not identify an energy
barrier, a formation rate, surface tension, or a critical group size
that separated growth from decay.

The link between market instability and a price crash depended on whether
investors could carry out their decisions. Traders first needed cash to buy
or shares to sell, after which buy and sell orders were matched between
traders. The dealer completed remaining orders only while its cash and share
limits allowed. These limits connect the model to accounts of how dealer
holdings affect
trading~\cite{garman1976market,stoll1978supply,ho1981optimal,braunmunzinger2018agent}.
Only the difference between completed trader buys and sells changed the
price~\cite{kyle1985continuous,farmer2005predictive,bouchaud2009how}.
Without dealer capacity, equal numbers of completed buys and sells left the
price fixed, whereas unlimited capacity allowed every order backed by trader
resources to complete. Runs from the same illustrated starting market
consequently reached a Sell majority with either a 0.40\% or a 25.8\% drawdown (Fig.~\ref{fig:dealer-panels}). At 18\% rule adoption with the original cash holdings, tenfold
trader shares produced an 83.5\% median drawdown against the unlimited
dealer, but limited capacity held the corresponding drawdowns to 3.92\% or 0.399\% (Fig.~\ref{fig:balance-sheet}). More shares allowed traders to place more sell orders, but selling
beyond what other traders bought lowered the price only while the dealer
could continue buying. A market dominated by sellers therefore did not
determine the severity of the crash on its own, because the price fell only
as far as someone was able to buy.

Completed trades changed the price, which could then influence later
decisions through the technical traders. At $f=0.175$, capacity $h=0.1$ raised majority probability by 38 percentage points relative to the unlimited dealer (Fig.~\ref{fig:reservoir}). The technical traders compared the price with a moving average and adjusted
the width of the band that triggered buying or selling. Keeping this
required departure fixed at its value before the rule change removed the
smaller dealer's advantage in the estimated probabilities, linking the advantage to the technical traders adjusting their bands (Supplementary Table~\ref{tab:revision-band}). The reference price
still changed, as did trader and dealer holdings, and the fixed departures
still differed across starting markets and capacities. Increasing capacity
also continued to produce both increases and decreases in majority
probability, so this adjustment did not explain the entire effect of dealer
capacity. Changing how strongly trades moved the price gave another limit on
the interpretation: larger dealers still produced larger drawdowns across a
tenfold range of price-impact coefficients, but drawdown magnitudes changed by a factor of seven to ten (Supplementary Table~\ref{tab:revision-impact}). Limiting completed trades could therefore
reduce drawdown while the resulting price changes made a Sell majority more
likely. The two outcomes measured how many traders decided to sell and how
far completed trades lowered the price. A Sell majority alone was
insufficient to identify a price crash.

The model differed from established agent-based markets in where traders
interacted and in what was changed. In the minority game, every trader reads
the same record of past outcomes, keeps a score for each of its strategies,
and uses the one that would have done best, so traders interact only through
the market outcome and the main control parameter is the number of
distinguishable histories relative to the number of
traders~\cite{challet1997minority,challet2000modeling}. In the Santa Fe
artificial stock market, traders learn forecasting rules from a shared price
and dividend~\cite{lebaron1999time}; adding technical rules to a trader's
repertoire was a dominant strategy that lowered wealth and raised volatility
for all~\cite{joshi1999technical}, and the rate at which traders revised
their rules settled at an equilibrium that left earnings below what slower
revision would have produced~\cite{joshi2002financial}. Models with
chartists and fundamentalists let traders move between the two groups at
rates set by recent profits or by the opinions of
others~\cite{lux1999scaling,brock1998heterogeneous,kirman1993ants}. Lattice
and cellular-automaton markets restricted influence to neighbors, but each
trader imitated its neighbors with some probability or responded to a spin
coupling at a set temperature, and the traders who used the price were
spread through the
grid~\cite{bornholdt2001ising,iori2002microsimulation,bartolozzi2004stochastic,qiu2007cellular,sornette2006importance,fan2009effect,atman2012influence}.
Wolfram's one-dimensional market and a later model built on Rule~110 used
deterministic rules and set the price from the numbers of buyers and
sellers~\cite{wolfram2002new,valsan2022on}. In each of these models a change
in behavior entered through learning or through a parameter applied to the
whole market, and its effect was read from the statistics of returns. Here
every trader kept a fixed rule, the intervention changed one entry of the
rule table for a chosen fraction of traders, and the outcome was compared
with a market that shared the same history. Only the technical traders on
the edge decided from the price, so the same rule change could be tested
with those traders active and held. Orders needed cash or shares and were
completed by other traders or by a dealer with limited holdings, so a Sell
majority and a price crash were measured separately. These choices gave up
the learning and strategic behavior of established models in exchange for a
rule space that can be enumerated and an intervention that can be repeated
exactly. The questions answered were correspondingly narrower: not whether
technical trading emerges or whether the market becomes efficient, but how
many traders had to change a rule before the population tipped, and whether
that tip required the price.

Several modeling choices limit how far these conclusions extend. Traders
followed assigned decision rules rather than learning a strategy or choosing
one to maximize returns. They updated together on a regular grid using one
original rule, and different combinations of neighboring decisions could
have the same sum and therefore produce the same response. Testing changes
to this lookup table did not cover all possible ways to respond to
neighboring traders. The technical traders followed one contrarian band
rule and occupied only the edge of the grid, so trend-following rules and
technical traders inside the grid were not examined. Each order requested one share at a common price, and
the dealer had no strategy for setting prices or managing risk. It could
replenish cash only by selling shares and replenish shares only by buying
them. The model included resource limits but no record of outstanding orders
at different prices, and its price response and time steps were not
calibrated to an observed market. Its simulated crashes therefore did not
predict the timing or magnitude of actual market crashes. Ten runs with the original decision rule in each condition, with the
technical traders either active or kept in Hold, remained below a Sell majority for 100,000 additional steps (Appendix~\ref{app:revision-controls}). This
extended the observed persistence of mixed decisions without establishing
that it would continue indefinitely, and neither the 1,000-step initial
period nor the final 500-step average established that the distribution of
market states had stopped changing. Updating traders at different times,
using other decision rules, or adding an order book would allow the same
question to be examined under different assumptions: how investor decisions
destabilize a market, and what determines whether that instability
accompanies a crash.
